\section{Background}

Federated learning (FL) trains a shared model across decentralised clients
without pooling raw data \cite{mcmahan2017communication}. It does not, by
itself, provide a privacy guarantee: model updates leak information about the
data that produced them, and reconstruction, membership-inference and
attribute-inference attacks all operate on shared gradients alone. Differential
privacy (DP) supplies the missing guarantee by clipping each per-sample gradient
to a norm $C$ and adding calibrated Gaussian noise \cite{abadi2016deep}, at a
cost in utility.

\vspace{10pt}

In the federated setting the privacy budget is consumed over many communication
rounds, which raises a scheduling question: given a total budget $\rho$, how
should it be divided among rounds? The intuition motivating adaptive schemes is
that gradients are large early in training and small later, so early rounds
should tolerate more noise. Several recent works formalise this as a learning
problem and train reinforcement-learning agents to allocate the budget
\cite{ChenAdaptivePrivacyMARL, WangDNORL}.

\vspace{10pt}

This thesis asks whether that intuition survives contact with the mechanism it
is built on, in a demanding application: federated thermal detection of small
unmanned aerial vehicles, where targets occupy a handful of pixels and the
detector has little accuracy to spare.

\section{Motivation}

Two observations motivate the work.

\vspace{10pt}

First, the adaptive-allocation intuition is stated in terms of \emph{gradient
magnitude}, but DP-SGD clips every per-sample gradient to a fixed norm before
noise is added. If clipping binds, and we measure it binding by a factor of
$27$ in our own setting and by a factor of $8$ in the source method's published
configuration, then the magnitude that the intuition refers to has been removed
from the mechanism before it can reach the aggregate. Whether anything remains
for an allocator to exploit is an empirical and theoretical question that the
literature has not settled.

\vspace{10pt}

Second, DP and Byzantine robustness are usually studied separately, and are
assumed to compose. They interact. DP noise is what a poisoning attacker can
hide behind, and averaging, the operation that a robust aggregator replaces, is
what suppresses DP noise. Neither interaction is captured by studying either
mechanism alone.

\section{Problem Statement}

Let $T$ communication rounds each consume $\rho_t$ of a client's total budget
$\rho$, with per-round noise multiplier $\sigma_t$ determined by $\rho_t$. The
allocation problem is to choose $\{\rho_t\}$ with $\sum_t \rho_t = \rho$ so as to
maximise final model utility.

\vspace{10pt}

The literature treats this as an open optimisation to be solved by a learned
policy. We ask instead whether it has a closed-form solution, and if so whether
the learned policies can improve on it. We then ask what the same architecture's
security properties are under adaptive model poisoning, and whether standard
Byzantine defences remain appropriate once DP is present.

\section{Research Objectives}

\begin{enumerate}[label=O\arabic*.]
  \item Characterise the optimum of the per-round privacy budget allocation
  problem analytically under per-sample clipping, rather than treating it as an
  open optimisation to be solved numerically.
  \item Test that characterisation empirically in a regime constructed so that
  accuracy is genuinely sensitive to the noise multiplier, using a statistical
  protocol precise enough to resolve the effect sizes at stake.
  \item Establish whether learned allocators, as published, can escape the
  resulting bound.
  \item Measure the honest update envelope of a differentially private detector
  and determine whether the camouflage premise of DP-exploiting poisoning
  attacks holds at detector scale.
  \item Evaluate detection-based and aggregation-based Byzantine defences in the
  presence of DP, including their cost to honest participants.
\end{enumerate}

\section{Contributions}

\begin{enumerate}
  \item \textbf{An optimality result.} Under per-sample clipping, uniform
  allocation minimises the noise variance accumulated into the global model, and
  every non-uniform schedule is strictly worse (Section~\ref{sec:theory}). The
  same inequality closes allocation across layers and across clients.

  \item \textbf{Empirical validation at scale.} Thirteen allocations from three
  unrelated families, paired with common random numbers, are ordered by a single
  scalar (Spearman $\rho=-0.767$, $p=0.0022$). No allocation beats uniform at
  any privacy budget tested (Section~\ref{sec:results-alloc}). The learned
  allocator itself, trained and evaluated on held-out data over paired trials,
  likewise fails to beat uniform (Section~\ref{sec:drl-practice}).

  \item \textbf{A configuration methodology for DP-FL detection.} We give an
  SNR model, $\text{eff.\ SNR}\propto \sqrt{B}/(\sigma\sqrt{D})$, that predicts
  where a detector's accuracy becomes noise-sensitive, and validate its
  predictions to within a factor of two (Section~\ref{sec:calibration}). We also
  identify a pitfall: adaptive optimisers such as Adam normalise the noise scale
  away entirely, rendering the privacy parameter inert.

  \item \textbf{A concentration result for DP poisoning.} High-dimensional
  concentration makes the honest update norm highly predictable ($0.10\%$
  relative standard deviation), inverting the premise of DP-camouflage attacks
  and capping attacker amplitude (Section~\ref{sec:p3-envelope}).

  \item \textbf{A Byzantine-resilience study under DP}, including the finding
  that classical robust aggregators are counterproductive in this regime
  (Section~\ref{sec:p3-results}).
\end{enumerate}

\section{Scope and Limitations}

Results are obtained on a clear-sky subset of the Anti-UAV600 thermal benchmark
\cite{zhao20233rd}, partitioned by recording session, with a frozen publicly
pretrained trunk and a 566,789-parameter privately trained head. Absolute
detection numbers describe this subset and should not be read as general
anti-UAV performance; the comparisons are relative and unaffected. The poisoning
study covers two attacks at up to $30\%$ malicious clients.
Section~\ref{sec:limitations} states what is and is not covered.

\section{Thesis Organisation}

Chapter~2 reviews differentially private federated learning, adaptive budget
allocation, Byzantine-robust aggregation, and poisoning attacks that exploit DP
noise, and identifies the gap addressed here. Chapter~3 formalises the system
model and proves the optimality of uniform allocation, together with its
extensions across layers and clients. Chapter~4 describes the dataset,
architecture and calibration procedure used to construct a regime in which the
allocation question is measurable at all, and states the statistical protocol.
Chapter~5 reports the allocation results and the Byzantine-resilience study.
Chapter~6 discusses implications, limitations and future work.
